% gwebluapdf.tex -- the LuaTeX PDF back end for GWEB.
%
% gwebmac's \dest/\link/\bookmark/\pdfURL are written for pdfTeX primitives
% that luatex does not provide; this file restates them with luatex's
% \pdfextension forms, so a web typeset with luatex (say, one that draws
% MetaPost figures through luamplib) gets the same blue cross-reference links
% and PDF outline (bookmarks) a pdfTeX run does. gwebmac \inputs it by itself
% once it detects a luatex run producing PDF, so no web has to ask for it; a
% Korean web gets it the same way, through kotexgweb.tex's own \input gwebmac.
%
% A PDF outline title must be a UTF-16BE string, so titles are converted by a
% small Lua helper, gweb_pdfutf16(s), which prints s as the body of a PDF
% literal string: the UTF-16BE byte-order mark followed by each code unit as a
% pair of octal-escaped bytes (surrogate-paired above the BMP). ASCII English
% titles pass through it unharmed; Korean titles need it.
%
% The helper is built inline. Two cautions about Lua inside \directlua: (1) the
% TeX-special characters % # _ & ^ ~ $ are made catcode 12 below (% is Lua's
% modulo operator, not a TeX comment; # is its length operator), and the
% backslash is produced as string.char(92) so none appears in the source; and
% (2) \directlua turns the source's newlines into spaces, so a Lua "--" comment
% would run to the end of the whole chunk -- hence there are none here.
\begingroup
\catcode`\%=12 \catcode`\#=12 \catcode`\_=12 \catcode`\&=12
\catcode`\^=12 \catcode`\~=12 \catcode`\$=12
\directlua{
  local function oct(v)
    local d1 = math.floor(v / 64); v = v % 64
    local d2 = math.floor(v / 8); local d3 = v % 8
    return tostring(d1) .. tostring(d2) .. tostring(d3)
  end
  function gweb_pdfutf16(s)
    local bs = string.char(92)
    local t = { bs .. oct(254) .. bs .. oct(255) }
    for _, c in utf8.codes(s) do
      if c >= 65536 then
        c = c - 65536
        local w1 = 55296 + math.floor(c / 1024)
        local w2 = 56320 + (c % 1024)
        t[#t + 1] = bs .. oct(math.floor(w1 / 256)) .. bs .. oct(w1 % 256)
        t[#t + 1] = bs .. oct(math.floor(w2 / 256)) .. bs .. oct(w2 % 256)
      else
        t[#t + 1] = bs .. oct(math.floor(c / 256)) .. bs .. oct(c % 256)
      end
    end
    tex.sprint(-2, table.concat(t))
  end
}
\endgroup
\def\dest#1{\pdfextension dest num #1 fith\relax}
\def\link#1#2{\pdfextension startlink attr{/Border[0 0 0]}goto num #1\relax
  \pdfextension literal{0 0 1 rg}#2\pdfextension literal{0 g}\pdfextension endlink}
\def\bookmark#1#2#3#4{\pdfextension outline goto num #2 count #3
  {\directlua{gweb_pdfutf16("\luaescapestring{#4}")}}}
% gwebmac's \pdfURL is likewise written for pdfTeX primitives; restate it (the
% link text may of course be Korean).
\def\pdfURL#1#2{%
  \pdfextension startlink attr{/Border[0 0 0]}
    user{/Subtype /Link /A << /S /URI /URI (#2) >>}\relax
  \pdfextension literal{0 0 1 rg}#1\pdfextension literal{0 g}%
  \pdfextension endlink}
% Open the bookmark (outline) pane when the document is first shown.
\pdfextension catalog{/PageMode /UseOutlines}
